% !TEX root = ../main.tex


\section{The Full Constraint List}

We use the domain and columns of Section 3.2. Each identity holds throughout \(H\). Table 11 identifies the rows on which it imposes a restriction. For limb width \(w\) (\(w=8\) when \(n\ge256\)), let \(m_{in}=\lceil k_{in}/w\rceil\) and \(m=\lceil k/w\rceil\). We set \(\nu=1\) for positive \(V_{min\Delta}\) and \(\nu=0\) otherwise.

\subsection{Scalar Checks}

The verifier first checks the public receipt against the following bounds and conventions.

\[0 \le c \le p^* \le d \le N-1, \qquad 1 \le V_{max} < 2^k, \qquad 0 \le V_{min\Delta} < 2^k,\]

\[
s \in \{0,1\}, \quad
V_{min\Delta} \le P < 2^k, \quad
0 \le J < 2^k, \quad
V_{min\Delta} = 0 \Rightarrow s = P = J = 0.
\]

\subsection{Clearing Constraints}

\textbf{Table 11.} The constraints of the clearing proof.

\begingroup
\small
\setlength{\tabcolsep}{3pt}

\begin{longtable}{@{}
  >{\raggedright\arraybackslash}p{2.8em}
  >{\raggedright\arraybackslash}p{0.18\linewidth}
  >{\raggedright\arraybackslash}p{\dimexpr0.66\linewidth-2.8em-6\tabcolsep\relax}
  >{\raggedright\arraybackslash}p{0.16\linewidth}
@{}}
\toprule
\# & Name & Identity & Active at \\
\midrule
\endhead
\bottomrule
\endlastfoot

F1 & Bid range
   & \(B - \sum_{l<m_{in}} 2^{wl}\ell^{B}_l = 0\)
   & All ticks \\

F2 & Ask range
   & \(A - \sum_{l<m_{in}} 2^{wl}\ell^{A}_l = 0\)
   & All ticks \\

F3 & Bid padding
   & \(Pad \cdot B = 0\)
   & Padding \\

F4 & Ask padding
   & \(Pad \cdot A = 0\)
   & Padding \\

F5 & Demand start
   & \(L_{n-1} \cdot (AccB - B) = 0\)
   & \(\omega^{n-1}\) \\

F6 & Supply start
   & \(L_0 \cdot (AccA - A) = 0\)
   & \(\omega^0\) \\

F7 & Demand step
   & \((X - \omega^{n-1})\,\bigl[AccB(X)\)\newline
     \(\quad {}- AccB(\omega X) - B(X)\bigr] = 0\)
   & All but the last \\

F8 & Supply step
   & \((X - \omega^{n-1})\,\bigl[AccA(\omega X)\)\newline
     \(\quad {}- AccA(X) - A(\omega X)\bigr] = 0\)
   & All but the last \\

F9 & Bid surplus
   & \(SurpB - (AccB - Min) = 0\)
   & All ticks \\

F10 & Ask surplus
    & \(SurpA - (AccA - Min) = 0\)
    & All ticks \\

F11 & Minimum selection
    & \(SurpA \cdot SurpB = 0\)
    & All ticks \\

F12 & \(SurpB \ge 0\)
    & \(SurpB - \sum_{l<m} 2^{wl}\ell^{sB}_l = 0\)
    & All ticks \\

F13 & \(SurpA \ge 0\)
    & \(SurpA - \sum_{l<m} 2^{wl}\ell^{sA}_l = 0\)
    & All ticks \\

F14 & Ceiling
    & \((V_{max} - Min) - \sum_{l<m} 2^{wl}\ell^{ceil}_l = 0\)
    & All ticks \\

F15 & Plateau, left end
    & \(L_c \cdot (Min - V_{max}) = 0\)
    & \(\omega^c\) \\

F16 & Plateau, right end
    & \(L_d \cdot (Min - V_{max}) = 0\)
    & \(\omega^d\) \\

F17 & Cliffs
    & \(Mask_C \cdot \bigl(V_{max} - Min\)\newline
      \(\quad {}- 1 - Slack\bigr) = 0\)
    & \(\omega^{c-1}\), \(\omega^{d+1}\) where they exist \\

F18 & Imbalance
    & \(\Delta - (SurpA + SurpB) = 0\)
    & All ticks \\

F19 & Imbalance floor
    & \(Mask_P \cdot \bigl[(\Delta - V_{min\Delta})\)\newline
      \(\quad {}- \sum_{l<m} 2^{wl}\ell^{flr}_l\bigr] = 0\)
    & Plateau \\

F20a & Pin, demand side
     & \(L_{p^*} \cdot (SurpB - (1-s)\,V_{min\Delta}) = 0\)
     & \(\omega^{p^*}\) \\

F20b & Pin, supply side
     & \(L_{p^*} \cdot (SurpA - s\,V_{min\Delta}) = 0\)
     & \(\omega^{p^*}\) \\

F21 & Tie-break
    & \(Mask_T \cdot \bigl(\Delta - V_{min\Delta}\)\newline
      \(\quad {}- 1 - Slack\bigr) = 0\)
    & \(\omega^{p^*-1}\) when \(Mask_T \ne 0\) \\

F22 & \(Slack \ge 0\)
    & \(Slack - \sum_{l<m} 2^{wl}\ell^{slk}_l = 0\)
    & All ticks \\

F23 & Table membership
    & The lookup identities of Section 3.4 over all looked-up values
    & All ticks \\

F24 & IO bid range
    & \(IOB - \sum_{l<m_{in}} 2^{wl}\ell^{IOB}_l = 0\)
    & All ticks \\

F25 & IO ask range
    & \(IOA - \sum_{l<m_{in}} 2^{wl}\ell^{IOA}_l = 0\)
    & All ticks \\

F26 & IO bid padding
    & \(Pad \cdot IOB = 0\)
    & Padding \\

F27 & IO ask padding
    & \(Pad \cdot IOA = 0\)
    & Padding \\

F28 & IO demand start
    & \(L_{n-1} \cdot (IOAccB - IOB) = 0\)
    & \(\omega^{n-1}\) \\

F29 & IO demand step
    & \((X - \omega^{n-1})\,\bigl[IOAccB(X)\)\newline
      \(\quad {}- IOAccB(\omega X) - IOB(X)\bigr] = 0\)
    & All but the last \\

F30 & IO supply start
    & \(L_0 \cdot (IOAccA - IOA) = 0\)
    & \(\omega^0\) \\

F31 & IO supply step
    & \((X - \omega^{n-1})\,\bigl[IOAccA(\omega X)\)\newline
      \(\quad {}- IOAccA(X) - IOA(\omega X)\bigr] = 0\)
    & All but the last \\

F32 & At-price long volume
    & \(\nu\,L_{p^*} \cdot \bigl((1-s)\,B\)\newline
      \(\quad {}+ s\,A - P\bigr) = 0\)
    & \(\omega^{p^*}\) when \(V_{min\Delta} > 0\) \\

F33 & Eligible IO volume
    & \(\nu\,L_{p^*} \cdot \bigl((1-s)\,IOAccA\)\newline
      \(\quad {}+ s\,IOAccB - J\bigr) = 0\)
    & \(\omega^{p^*}\) when \(V_{min\Delta} > 0\) \\

\end{longtable}
\endgroup

F20 is a pair, so the list has 33 identities besides the lookup. The looked-up values are the limbs of the nine range-checked quantities (\(B\), \(A\), \(IOB\), \(IOA\), \(SurpB\), \(SurpA\), \(V_{max} - Min\), \(\Delta - V_{min\Delta}\), \(Slack\)) and, for a width that is not a multiple of \(w\), the scaled top limb of each. Their number is \(E = 4 m_{in} + 5m\), plus 4 if \(w \nmid k_{in}\), plus 5 if \(w \nmid k\). F23 consists of \(\lceil E/2 \rceil\) pair identities, the table-helper identity, and the running-sum identity.

\textbf{Column count.} The first round commits to 13 value columns, \(4m_{in}+5m\) limb columns, and \(Mult\). The second adds \(\lceil E/2\rceil+2\) helper columns. At \(n=1024\), \(k_{in}=32\), and \(k=42\), we have 46 limb columns and \(E=51\) lookup inputs. The 26 pair helpers bring the total, including the other helpers, to 88 columns.

\textbf{Degrees and openings.} The pair identities have degree at most \(3n - 3\) before masking and every other identity at most \(2n - 2\). The quotient has degree at most \(2n + 3\) with masking, and the implementation requests a reference string of degree \(2n + 8\). \(AccB\), \(AccA\), \(IOAccB\), \(IOAccA\), \(A\), \(IOA\), and \(S\) are opened at \(\zeta\) and \(\omega\zeta\); every other column at \(\zeta\).

\subsection{No-Cross Constraints}

The no-cross proof uses only ordinary bids and asks, their depth columns, the limbs of \(B\) and \(A\), and the corresponding lookup columns. We impose F1--F8, apply F23 to those limbs, and require

\[AccA(X) \cdot AccB(X) = 0.\]

\subsection{Update-Validity Constraints}

For update validity, let \(\kappa\) generate the domain \(K\) of size \(M\), with \(\kappa^{M/n}=\omega\). The public polynomial \(P_\eta\) takes value \(\eta^j\) on the first \(\mu\) rows, one per logged step, and zero elsewhere. On row \(iM/n\), \(X_H\) equals \(\omega^i\) and \(S_H\) equals one. Both vanish on the other rows. The range table is also public. The witness includes \(Qty\), \(Tick\), \(W\), quantity limbs and their lookup columns. The prover commits to \(G\), \(U_H\), and \(S_2\) after receiving \(\beta\).

\textbf{Table 12.} The constraints of the update-validity proof, each on all of \(K\).

\begin{longtable}[]{@{}lll@{}}
\toprule\noalign{}
\# & Name & Identity \\
\midrule\noalign{}
\endhead
\bottomrule\noalign{}
\endlastfoot
U1 & Quantity range & \(Qty - \sum_{l<m_{in}} 2^{wl} \ell^{Qty}_l = 0\), with every limb in the table \\
U2 & Step side & \(G(X)\,(\beta - Tick(X)) - P_\eta(X)\,Qty(X) = 0\) \\
U3 & Tick side & \(U_H(X)\,(\beta - X_H(X)) - S_H(X)\,W(X) = 0\) \\
U4 & Running sum & \(S_2(\kappa X) - S_2(X) - (G(X) - U_H(X)) = 0\) \\
\end{longtable}

For 32-bit quantities the proof has 15 columns: \(Qty\), \(Tick\), \(W\), four limbs, the multiplicities, two pair helpers, the table helper, the range running sum, \(G\), \(U_H\), and \(S_2\). With the quotient and \(\Phi_{mask}\) it holds 17 commitments.


